\newpage
\section{Change of Groups}

Suppose that $G$ is a group acting on an abelian group $A$,
and suppose that there is a group homomorphism $\phi: G' \to G$.
Then $\phi$ induces a $G'$-action on $A$.

\begin{question}
Any connection between $H^i(G, A)$ and $H^i(G', A)$?
\end{question}

For $i = 0$, we have $H^0(G, A) = A^G$ and $H^0(G', A) = A^{G'}$,
and we have an inclusion $A^G \hookrightarrow A^{G'}$.

\subsection{Change of Rings}

We consider a more general setting as before.
Let $\alpha: R \to S$ be a ring homomorphism and they don't have to be commutative.
We are given $S$-modules $P, M$ and an $R$-module $N$.
We have an $R$-linear map $h: M \to N$.

\begin{theorem}
    Given the above data,
    there is a natural abelian group homomorphism $h_*: \ext_S^i(P, M) \to \ext_R^i(P, N)$ for all $i \ge 0$.
    Moreover, if $R$ is commutative, then $h_*$ is an $R$-linear map.
\end{theorem}
\begin{proof}
    Let $M \to I^\bullet$ be an injective resolution of $M$ as an $S$-module.
    Let $I^\bullet \to J^\bullet$ be an injective resolution of $I^\bullet$ as an $R$-module.
    Then we have canonically induced maps:
    \[
    \begin{tikzcd}
        \Hom_S(P, M) \arrow{r} & \Hom_R(P, I^\bullet)\arrow{d}{\alpha^*}\\
        & \Hom_R(P, I^\bullet)\arrow{d}\\
        & \Hom_R(P, J^\bullet)
    \end{tikzcd}
    \]
    hence we have $H^i(\Hom_S(P, I^\bullet))(=\ext_S^i(P, M)) \to H^i(\Hom_R(P, J^\bullet))(=\ext_R^i(P, M))$.
    With $h: M \to N$, a map $\ext^i_R(P, M) \to \ext^i_R(P, N)$ is induced.
    Thus we have a map $h_*: \ext_S^i(P, M) \to \ext_R^i(P, N)$ by composition.
\end{proof}

We note that $h_*$ is a compusition of $\Hom_S(P, M) \hookrightarrow \Hom_R(P, M)$
and $h\circ (-): \Hom_R(P, M) \to \Hom_R(P, N)$ when $i=0$.

\subsection{Restriction Maps}
\begin{proposition}
    Let $\varphi: G' \to G$ be a group homomorphism and $A$ be a $G$-module.
    Then we have a natural map $\varphi^*: H^i(G, A) \to H^i(G', A)$ for all $i \ge 0$.
\end{proposition}
\begin{proof}
    Take $R = \Z[G']$, $S = \Z[G]$, $P = \Z$, $M = A$, and $N = B$.
    Then we have a natural map $\varphi^*: \ext^i_{\Z[G]}(\Z, A) \to \ext^i_{\Z[G']}(\Z, B)$ for all $i \ge 0$.
\end{proof}

We note that when $i=0$, we have $A^G \hookrightarrow A^{G'} \overset{h_*}{\hookrightarrow} B^{G'}$.

\begin{example}[Restriction Maps]\noindent
    \begin{enumerate}
        \item Take $A = B$ and $h = \mathrm{id}_A$.
        Then we have a natural map $\Res_{G'}^G: H^i(G, A) \to H^i(G', A)$ for all $i \ge 0$.
        As we have seen, when $i=0$, we have $A^G \hookrightarrow A^{G'}$.
        \item Take $G'$ to be a subgroup $H$ of $G$.
        Then we get a restriction map $\Res_H^G: H^i(G, -) \to H^i(H, -)$ on $G$-modules.
    \end{enumerate}
\end{example}

\begin{example}[Topological Interpretation of Restriction Maps]
    
Let $H$ be a subgroup of $G$.
Recall from Section 1 that $H^i(G, \Z) \cong H^i_{sing}(BG, \Z)$.
So the restriction map $\Res_H^G$ should be a map
\[
    H^*_{sing}(BG, \Z) \to H^*_{sing}(BH, \Z),
\]
and a natural way to produce such a map in topology is to pull back along a continuous map.

\begin{question}
    Is there a continuous map $BH \to BG$?
\end{question}

The tool is the universal property of $BG$ which classifies principal $G$-bundles.
That is, for a (paracompact) space $X$,
\[
    [X, BG] \xrightarrow{\;\sim\;} \{\text{principal $G$-bundles over } X\}/\cong,
    \qquad [f] \mapsto f^*EG.
\]
Hence to construct a map $BH \to BG$ (up to homotopy),
it is enough to construct a principal $G$-bundle over $BH$.
Given a principal $H$-bundle $EH \to BH$, we want to enlarge its structure group from $H$ to $G$.
We can do this by taking the associated bundle $EH \times^H G$ with fiber $G$:
$EH \overset{H}{\times} G \to BH$, $[e, g] \mapsto [e]$.

By the universal property, there is a continuous map $\iota: BH \to BG$,
unique up to homotopy, such that $\iota^* EG \cong EH \overset{H}{\times} G$ as principal $G$-bundles.
Here are some remarks on this construction.
    \begin{enumerate}
        \item A concrete model: $G$ acts freely on the contractible space $EG$, hence so does $H$.
        So $EG$ is also a model for $EH$, and $BH = EG/H$.
        In this model $\iota$ is simply the quotient map
        \[
            BH = EG/H \longrightarrow EG/G = BG,
        \]
        which is a covering map with fiber $G/H$ (of degree $[G:H]$ when it is finite).
        \item The map $\iota$ induces $\iota^*: H^*_{sing}(BG, \Z) \to H^*_{sing}(BH, \Z)$,
        and under the identifications $H^*(G, \Z) \cong H^*_{sing}(BG, \Z)$, $H^*(H, \Z) \cong H^*_{sing}(BH, \Z)$
        this is the restriction map $\Res_H^G$.
        Indeed, on $\pi_1$ the map $\iota$ is the inclusion $H \hookrightarrow G$,
        and $\Res_H^G$ is the map on cohomology induced by this group homomorphism.
    \end{enumerate}
\end{example}

\subsection{Inflication Maps}
Suppose that $H$ is a normal subgroup of $G$,
so we can consider the quotient map $\pi: G \to G/H$.